\section*{Demostració que \(z_1^n + z_2^n\) és real}

Siguin $z_1, z_2 \in \mathbb{C}$ les arrels del polinomi
$x^2 + ax + b$
amb $a, b \in \mathbb{R}$. Com que els coeficients són reals, hem vist a classe que si un nombre complex és arrel també ho serà el seu conjugat. Per tant, sense pèrdua de generalitat (es faria de manera anàloga amb $z_1 = \overline{z_2}$)
\begin{equation*}
z_2 = \overline{z_1}.
\end{equation*}
Per a qualsevol $n\in \mathbb{Z^+}$ considerem la suma
\begin{equation*}
z_1^n + z_2^n = z_1^n + \overline{z_1}^n.
\end{equation*}

Com que $\overline{z_1}^n = \overline{z_1^n}$


Llavors,
\begin{equation*}
z_1^n + z_2^n = z_1^n + \overline{z_1^n}.
\end{equation*}

Però per a qualsevol nombre complex $(w)$, es té que $(w + \overline{w} = 2,\operatorname{Re}(w) \in \mathbb{R})$. En particular, prenent $(w = z_1^n)$, obtenim
\begin{equation*}
z_1^n + z_2^n = 2\,\operatorname{Re}(z_1^n) \in \mathbb{R}.
\end{equation*}

\subsection*{Cas concret}
Per al polinomi $(x^2 - 2x + 2)$, les arrels són
\begin{equation*}
z_1 = 1 + i, \qquad z_2 = 1 - i = \overline{z_1}.
\end{equation*}
que ho hem resolt mitjançant la formula 

\begin{equation*}
x = \frac{2 \pm \sqrt{4^2 -4*1*2}}{2}
\end{equation*}
com tenim una arrel negativa, hem de utilitzar complexos
\newline

\begin{equation*}
x = \frac{2 + 2i}{2} = 1 + i
\end{equation*}
\newline
\begin{equation*}
x = \frac{2 - 2i}{2} = 1 - i
\end{equation*}
\newline
que és real per a tot n en els enters, com es comprova directamen, farem la distinció parell de senar (per veure el valor de l'expressió en funció de n):
\subsection*{Parell}

n = 2k \(k \in \mathbb{Z}^+\)
aleshores tenim $(a_1 + ib_1)^n = $

\subsection*{Senar}

n = 2k + 1 \(k \in \mathbb{Z}^+\)
en aquest cas, la part real es multiplica per 2 mentre que la part imaginaria s'anula